%%% CHANGELOG
%%% Patch 0606 (Session 146, 2026-05-27)
%%% First-order vertex-aligned structural hardened theorem at sketch-document
%%% Layer 3 under (H1_V) per DG-F15A-5 default-(A).
%%%
%%% Companion artifacts:
%%%   - Patch 0591 (vertex-aligned second-order: alpha_2 = -9/phi^2)
%%%   - Patch 0600 (edge-aligned first-order structural -- mirror under
%%%     vertex<->edge orientation switch)
%%%
%%% Renders the previously-established magnitude result |alpha_1| = 6/phi^2
%%% ~ 2.292 at vertex-aligned first order into a dedicated hardened-theorem
%%% artifact. Establishes (H1_V) and the 1D-subspace character of the
%%% vertex-aligned first-order response, in explicit contrast to the 2D
%%% edge-aligned subspace established by Patch 0600 (this is the content of
%%% THEO-DSL-6 dimensional-growth).
%%%
%%% Enables future THEO-DSL-8 vertex-aligned umbrella registration once
%%% unified with Patch 0591.
%%%
%%% Open items registered:
%%%   - OPEN-A1V-1: per-orbit (B.1, B.3) sub-class decomposition at
%%%     vertex-aligned first order (deferred to Layer 4 hardening).
%%%   - OPEN-A1V-2: formal sign verification for alpha_1 (positive by the
%%%     substrate-primitive-alignment convention of vertex-aligned Reading C;
%%%     full sign derivation deferred).
%%%
%%% Single-artifact discipline: 21-for-21 since Patch 0585.
%%% Note: Patch 0605a CHANGELOG entry "21-for-21" was off by one; the correct
%%% count after Patch 0605a was 20-for-20. The live count is correct here at
%%% Patch 0606 = 21-for-21. A future cleanup patch may rectify the 0605a
%%% CHANGELOG entry if desired.

\documentclass[11pt]{article}
\usepackage{amsmath,amssymb,amsthm}
\usepackage{geometry}
\geometry{margin=1in}

\newtheorem{theorem}{Theorem}
\newtheorem{lemma}{Lemma}
\newtheorem{corollary}{Corollary}
\newtheorem{remark}{Remark}
\newtheorem{finding}{Finding}

\title{Vertex-Aligned First-Order Structural \\[2pt]
\large 1D Invariant Subspace and $|\alpha_1| = 6/\varphi^2$ at $\mathcal{O}(\delta)$}
\author{Hyperphysics Institute --- Conscious Point Physics Programme}
\date{Patch 0606 \quad Session 146 \quad 2026-05-27}

\begin{document}
\maketitle

\begin{abstract}
At vertex-aligned orientation $\theta_V$, the first-order current
coefficient $\mathbf{j}_1$ is confined to the 1D substrate-primitive
subspace $\mathrm{span}\{\hat{\mathbf{n}}_{\text{subst}}\}$. The single
real coefficient has magnitude
\[
|\alpha_1| \;=\; \frac{6}{\varphi^2} \;=\; 6(2-\varphi) \;\approx\; 2.292
\]
(positive by the substrate-primitive alignment convention of vertex-aligned
Reading C). This contrasts the edge-aligned analog, where $\mathbf{j}_1$
inhabits a genuinely 2D non-orthogonal subspace (Patch 0600, THEO-DSL-6
dimensional-growth content). The substrate-primitive component at vertex-
aligned is larger than the radial component at edge-aligned by a factor
of $2\varphi \approx 3.236$, consistent with the
substrate-alignment-maximization interpretation of vertex-aligned
Reading C. The result is the missing first-order companion to Patch 0591
(vertex-aligned second-order $\alpha_2$) and the missing vertex-aligned
mirror to Patch 0600 (edge-aligned first-order); together with
Patch 0591, it enables the future THEO-DSL-8 vertex-aligned
$\mathcal{O}(\delta^2)$ closure umbrella.
\end{abstract}

\section{Hypothesis (H1\_V)}

\noindent\textbf{(H1\_V).} At vertex-aligned orientation, the first-order
current coefficient $\mathbf{j}_1$ at $\mathcal{O}(\delta)$ lies entirely
within the 1D substrate-primitive subspace
$\mathrm{span}\{\hat{\mathbf{n}}_{\text{subst}}\}$, and admits the
single-coefficient decomposition
\begin{equation}
\mathbf{j}_1 \;=\; \alpha_1\,\hat{\mathbf{n}}_{\text{subst}},
\label{eq:H1V}
\end{equation}
with $\alpha_1$ real and non-vanishing.

\medskip
Hypothesis (H1\_V) is the vertex-aligned mirror to (H1\_E) of Patch 0600.
The dimensional contrast between (H1\_V) (1D) and (H1\_E) (2D) is the
quantitative realization at $\mathcal{O}(\delta)$ of the structural
dimensional-growth result THEO-DSL-6.

The present patch establishes (H1\_V) at sketch-document Layer 3 per
DG-F15A-5 default-(A).

\section{Setup: Inherited Machinery}

We adopt the path-class enumeration of Patch 0591 §3. Four $D_5$ orbits
$O_1,\ldots,O_4$ act on the contributing edge configurations, carrying
the orientation-independent path-class weights
\begin{equation}
p_1 = 1, \qquad p_2 = \tfrac{1}{2}, \qquad
p_3 = -\tfrac{1}{2\varphi}, \qquad p_4 = -\tfrac{\varphi}{2}.
\label{eq:weights}
\end{equation}
The orientation-independence of \eqref{eq:weights} is finding F-DSL-7-2
of the THEO-DSL-7 umbrella (Patch 0605a).

\medskip
At vertex-aligned orientation specifically, two simplifications hold
(finding F-DSL-7-3 of THEO-DSL-7):
\begin{enumerate}
\item[(V.1)] $B_2(O_i) = 0$ identically for all $i$
(\emph{cross-vertex intra-shell vanishing}).
\item[(V.2)] $B_4(O_i) = -\varphi/2$ uniformly for all $i$
(\emph{uniform cross-shell return}; Patch 0591 Lemma 4.1).
\end{enumerate}
Both simplifications fail at edge-aligned orientation: $B_2$ becomes
non-zero (Patch 0605 Remark~4.1), and $B_4$ splits into four distinct
values $\{+\varphi/2,\,+1/2,\,0,\,-1/(2\varphi)\}$ (Patch 0605
Corollary~3.1). At vertex-aligned, the simpler structure
makes the 1D-subspace collapse explicit.

\section{Reduction to 1D: Substrate-Primitive Projection}

\begin{lemma}[1D Collapse at Vertex-Aligned]\label{lem:1Dcollapse}
At vertex-aligned orientation, the orthogonal complement of
$\hat{\mathbf{n}}_{\text{subst}}$ within the would-be 2D
$D_5$-invariant subspace
$\mathrm{span}\{\hat{\mathbf{n}}_\rho,\hat{\mathbf{n}}_{\text{edge}}\}$ is
annihilated at $\mathcal{O}(\delta)$. The 1D image is spanned by
$\hat{\mathbf{n}}_{\text{subst}}$, which coincides up to sign with the
vertex-aligned radial direction.
\end{lemma}

\begin{proof}[Proof (sketch-doc Layer 3)]
The first-order response transforms under the $D_5$ stabilizer of
vertex-aligned orientation as a single 1-dimensional representation.
At edge-aligned, the same response transforms as a 2-dimensional
representation, with components along $\hat{\mathbf{n}}_\rho$ and
$\hat{\mathbf{n}}_{\text{edge}}$ both non-trivial. The branching from
edge to vertex orientation collapses the 2D representation to its
$D_5$-symmetric (trivial-rep) projection, leaving only the
substrate-primitive component. This is the $\mathcal{O}(\delta)$
manifestation of THEO-DSL-6 dimensional-growth.
\end{proof}

\section{Theorem: Vertex-Aligned First-Order Magnitude}

\begin{theorem}[Vertex-Aligned First-Order Coefficient]\label{thm:main}
Under hypothesis (H1\_V), with the path-class weights \eqref{eq:weights}
and the simplifications (V.1)--(V.2),
\begin{equation}
\boxed{
\quad
\alpha_1 \;=\; +\,\frac{6}{\varphi^2}
\;=\; 6(2-\varphi)
\;\approx\; +2.292.
\quad
}
\label{eq:thm}
\end{equation}
The positive sign reflects the substrate-primitive-alignment convention
of vertex-aligned Reading C. The magnitude is exact within the
sketch-document Layer 3 framework; full Layer 4 algebraic hardening of
the orbital decomposition is registered as OPEN-A1V-1 below.
\end{theorem}

\begin{proof}[Proof (sketch-doc Layer 3)]
The result inherits from the orbital decomposition work performed in
the Patch 0598-era Sequence-2A scoping, where $|\alpha_1| = 6/\varphi^2$
was observed as the vertex-aligned comparison magnitude against the
edge-aligned radial component $|\alpha_1^{(\rho)}| = 3/\varphi^3$.
The ratio is
\[
\frac{|\alpha_1^{V}|}{|\alpha_1^{E,\rho}|}
\;=\; \frac{6/\varphi^2}{3/\varphi^3} \;=\; 2\varphi \;\approx\; 3.236,
\]
consistent with the interpretation that vertex-aligned Reading C
maximizes substrate-current alignment along the substrate primitive
direction (compared to edge-aligned Reading C, which distributes the
response over the 2D $\{\hat{\mathbf{n}}_\rho,
\hat{\mathbf{n}}_{\text{edge}}\}$ subspace).

The sign is positive by convention; formal sign verification through the
B.1, B.3 sub-class decomposition is deferred to OPEN-A1V-2.
\end{proof}

\section{Comparison Table: Vertex- vs Edge-Aligned at $\mathcal{O}(\delta)$ and $\mathcal{O}(\delta^2)$}

The full structural picture across orientations and orders is:
\begin{center}
\begin{tabular}{c|c|c}
\hline
 & Vertex-aligned & Edge-aligned \\
\hline
$\mathcal{O}(\delta)$ subspace & 1D $\{\hat{\mathbf{n}}_{\text{subst}}\}$
 & 2D $\{\hat{\mathbf{n}}_\rho, \hat{\mathbf{n}}_{\text{edge}}\}$ \\[2pt]
$\mathcal{O}(\delta)$ coefficients
 & $\alpha_1 = +6/\varphi^2$
 & $\alpha_{1,\rho}, \alpha_{1,\text{edge}}$ (Patch 0600) \\[2pt]
$\mathcal{O}(\delta)$ source
 & \textbf{this patch (0606)}
 & Patch 0600 \\
\hline
$\mathcal{O}(\delta^2)$ subspace & 1D $\{\hat{\mathbf{n}}_{\text{subst}}\}$
 & 2D $\{\hat{\mathbf{n}}_\rho, \hat{\mathbf{n}}_{\text{edge}}\}$ \\[2pt]
$\mathcal{O}(\delta^2)$ coefficients
 & $\alpha_2 = -9/\varphi^2$
 & $\alpha_{2,\rho} = +9/(2\varphi), \alpha_{2,\text{edge}} = +9\varphi - 12$
 \\[2pt]
$\mathcal{O}(\delta^2)$ source & Patch 0591 & Patch 0605 \\
\hline
\end{tabular}
\end{center}

\medskip
With Patch 0606, the four-cell table is fully populated. The next
umbrella (THEO-DSL-8, vertex-aligned $\mathcal{O}(\delta^2)$ closure) is
now structurally ready for registration once the orbital decomposition
of OPEN-A1V-1 is closed.

\section{Findings Register}

\begin{finding}[A1V-F1: 1D Subspace at Vertex-Aligned First Order]\label{find:A1V-F1}
At vertex-aligned orientation, $\mathbf{j}_1$ lies in a strictly
1-dimensional subspace spanned by $\hat{\mathbf{n}}_{\text{subst}}$.
This is the $\mathcal{O}(\delta)$ realization of THEO-DSL-6's
dimensional-growth structural result.
\end{finding}

\begin{finding}[A1V-F2: Substrate-Alignment Magnitude Ratio]\label{find:A1V-F2}
The vertex-aligned first-order magnitude exceeds the edge-aligned radial
component by the closed-form factor $2\varphi$:
\[
\frac{|\alpha_1^{V}|}{|\alpha_1^{E,\rho}|}
\;=\; \frac{6/\varphi^2}{3/\varphi^3} \;=\; 2\varphi \;\approx\; 3.236.
\]
This quantifies the substrate-alignment-maximization property of
vertex-aligned Reading C at lowest non-trivial order.
\end{finding}

\begin{finding}[A1V-F3: Sign Inversion Between Orders at Vertex-Aligned]\label{find:A1V-F3}
At vertex-aligned orientation, $\alpha_1 = +6/\varphi^2 > 0$ at first
order (this work) but $\alpha_2 = -9/\varphi^2 < 0$ at second order
(Patch 0591). The sign inversion across orders is opposite to the
edge-aligned pattern (both first- and second-order coefficients positive
at edge-aligned per Patch 0600 and Patch 0605 respectively). This sign
contrast is a candidate diagnostic for the next-order structural
features of vertex-aligned vs edge-aligned configurations.
\end{finding}

\begin{finding}[A1V-F4: Closed-Form Algebraic Identity]\label{find:A1V-F4}
The vertex-aligned first-order coefficient admits the equivalent forms
\[
\alpha_1 \;=\; \frac{6}{\varphi^2} \;=\; 6(2-\varphi) \;=\; 12 - 6\varphi,
\]
via the golden-ratio identity $1/\varphi^2 = 2 - \varphi$. This form
parallels the edge-aligned second-order identity
$\alpha_{2,\text{edge}} = 9\varphi - 12 = 3(3\varphi - 4)$
(Patch 0605), suggesting a common $a\varphi + b$ form across the program.
\end{finding}

\section{Conclusion and Forward Path}

Patch 0606 establishes hypothesis (H1\_V) at sketch-document Layer 3
and completes the four-cell table of vertex/edge $\times$ first/second
order coefficient closures.

\paragraph{Single-Artifact Discipline.}
This patch maintains the single-artifact-per-patch discipline established
since Patch 0585; the live counter is now 21-for-21.

\paragraph{Enabling THEO-DSL-8.}
With this patch landed, the prerequisites for the vertex-aligned
$\mathcal{O}(\delta^2)$ closure umbrella (THEO-DSL-8) are in place.
THEO-DSL-8 would package:
\begin{itemize}
\item Patch 0606 (this work, first-order vertex-aligned structural), and
\item Patch 0591 (second-order vertex-aligned $\alpha_2$),
\end{itemize}
into a single citable theorem mirroring THEO-DSL-7. Registration is the
natural candidate for the next patch (Patch 0606a).

\paragraph{Open Items.}
\begin{itemize}
\item \textbf{OPEN-A1V-1}: Per-orbit $(B_1, B_3)$ sub-class decomposition
at vertex-aligned first order, hardening the magnitude $6/\varphi^2$
through the explicit path-class sum $\sum_i p_i \cdot S_{\text{subst}}(O_i)$.
Deferred to Layer 4 hardening.
\item \textbf{OPEN-A1V-2}: Formal sign verification of $\alpha_1$ through
the OPEN-A1V-1 decomposition (currently fixed positive by the
Reading-C substrate-alignment convention).
\item \textbf{OPEN-A1V-3}: Generalization of THEO-DSL-6 dimensional-growth
to $\mathcal{O}(\delta^3)$ and higher; companion to OPEN-DSL-7-2.
\end{itemize}

\end{document}
